% GENERATED FILE. Edit canonical lessons, then run npm run generate:books.
% canonical-source-hash: fnv1a64:fb46f13f6e530c64
% canonical-lessons: BN-C272-obhigyota, BN-W272-letter-o-indep, BN-C272-bhalo-lagchhe, BN-C272-hote-chai, BN-R272-feelings-and-hopes

\chapter{Feelings and Hopes}
\label{ch:bn-feelings-and-hopes}
\hlchaptermodality{272}

\begin{chapteropening}
\textbf{By the end of this chapter:} \emph{I can say how I feel with āmār (āmār bhālo lāgchhe, āmār bhôy kôrchhe), say what I want to become (ḍāktār hôte chāi) and what I wish, and read and write \bn{অ} in \bn{অভিজ্ঞতা}.}

Complete BN-R272-feelings-and-hopes: say how you feel, what you want to become and what you wish.
\end{chapteropening}

\section[\texorpdfstring{\bn{অভিজ্ঞতা} (ôbhigyôtā)}{অভিজ্ঞতা (ôbhigyôtā)}]{\bn{অভিজ্ঞতা} (ôbhigyôtā) — experience}

\label{lesson:BN-C272-obhigyota}



\begin{quote}
Before the new one: say \textquotedblleft{}My train has left\textquotedblright{}. (\emph{āmār ṭren chheṛe gechhe}.) Then \textbf{\bn{জ্ঞান}}, \textquotedblleft{}knowledge\textquotedblright{}: its \textbf{\bn{জ্ঞ}} is in this word too.
\end{quote}

\subsection*{You'll want to know: \bn{অভিজ্ঞতা}}
\textbf{\bn{অভিজ্ঞতা}} — \emph{ôbhigyôtā} — \textquotedblleft{}experience; an experience\textquotedblright{}.

\begin{quote}
\textbf{\bn{গ্রামে} \bn{আমার} \bn{একটা} \bn{নতুন} \bn{অভিজ্ঞতা} \bn{হল।}}
\end{quote}

\begin{quote}
\emph{grāme āmār ekṭā nôtun ôbhigyôtā hôlô.} — \textquotedblleft{}I had a new experience in the village.\textquotedblright{}
\end{quote}

An experience \textquotedblleft{}happens to\textquotedblright{} you, as a thought comes to you: \textbf{\bn{আমার} … \bn{অভিজ্ঞতা} \bn{হল}}. Experience you have gathered \textquotedblleft{}is there\textquotedblright{}: \textbf{\bn{এই} \bn{কাজে} \bn{তার} \bn{অভিজ্ঞতা} \bn{আছে}} (\emph{ei kāje tār ôbhigyôtā āchhe}), \textquotedblleft{}he has experience in this work\textquotedblright{}.

\begin{cousinweb}
\textbf{\bn{অভিজ্ঞতা}} (\emph{ôbhigyôtā}) is a tatsama: Sanskrit \emph{abhijñatā}, which Monier-Williams gives as \textquotedblleft{}the knowledge of\textquotedblright{}. It is \emph{abhi-}, \textquotedblleft{}towards\textquotedblright{}, with \emph{jña}, \textquotedblleft{}knowing\textquotedblright{} (the \textbf{\bn{জ্ঞ}} of \textbf{\bn{জ্ঞান}}, said \emph{gy}), and \emph{-tā}, which makes an abstract noun. Bengali uses it for knowledge got by living through something; that sense, \textquotedblleft{}experience\textquotedblright{}, is Bengali use, not one Monier-Williams gives.
\end{cousinweb}

\subsection*{Your turn}
Say these aloud:

\begin{itemize}
\raggedright
  \item \emph{ôbhigyôtā}
  \item \textquotedblleft{}I had a new experience in the village\textquotedblright{} — \emph{grāme āmār ekṭā nôtun ôbhigyôtā hôlô}
  \item \textquotedblleft{}He has experience in this work\textquotedblright{} — \emph{ei kāje tār ôbhigyôtā āchhe}
\end{itemize}

\begin{tcolorbox}[breakable,colback=teal!4,colframe=teal!35!black,title={Before you move on}]
What does \textbf{\bn{অভিজ্ঞতা}} mean? (\textquotedblleft{}Experience\textquotedblright{}.) \textquotedblleft{}I had an experience\textquotedblright{}? (\textbf{\bn{আমার} \bn{একটা} \bn{অভিজ্ঞতা} \bn{হল}}.)
\end{tcolorbox}
\section[\texorpdfstring{\bn{অ} (ô)}{অ (ô)}]{\bn{অ} (ô) — the letter at the start of \bn{অভিজ্ঞতা}}

\label{lesson:BN-W272-letter-o-indep}



\begin{quote}
Before the new one: say \textbf{\bn{অভিজ্ঞতা}}. Then read \textbf{\bn{ক}} (\emph{kô}): its vowel, the \emph{ô} every consonant carries, is the vowel of this letter.
\end{quote}

\subsection*{Script you'll notice: \bn{অ}}
\textbf{\bn{অ}} — \emph{ô}, the vowel inside every bare consonant: \textbf{\bn{ক}} is \emph{kô}. It has no sign, because a consonant carries it already; at the start of a word, with nothing to ride on, it is the letter \textbf{\bn{অ}}: \textbf{\bn{অভিজ্ঞতা}}. Now the name Arun can be written: \textbf{\bn{অরুণ}} (\emph{ôrun}).

In the print of this book, \textbf{\bn{অ}} has a headline and an upright stem on the right, the full height of the letter. Left of the stem is a large rounded bowl; its left arm ends in a slanted tip at the upper left. The bowl's right side is curled over inside the top of the bowl and ends there in a solid round knob, and a short curve joins the knob to the headline above it. A slanting leg joins the bottom right of the bowl to the foot of the stem.

\textbf{\bn{আ}}, which you write, is \textbf{\bn{অ}} with one more upright stroke after it, standing from the headline: look at the two side by side.

\subsection*{Writing: \bn{অ} — copy the printed shape}
Put your pen beside \textbf{\bn{অ}} and copy the shape you can see, slowly and large. Check the copy against the print: the bowl and its tip, the knob and the short curve to the headline, the leg to the stem.

\begin{quote}
This book does not yet tell you where to start \bn{অ} or which way the pen goes. That is not written down here until it can be written down with a source. Copying the printed shape needs no such source, and it is how the shape becomes yours in the meantime.
\end{quote}

\subsection*{Your turn}
\begin{itemize}
\raggedright
  \item \emph{Find:} \bn{অ} at the front of \textbf{\bn{অভিজ্ঞতা}}
  \item \emph{Write it:} \bn{অ}, then \bn{আ}, from the print
  \item \emph{Write it:} \textbf{\bn{অভিজ্ঞতা}}
  \item \emph{Write it:} \textbf{\bn{অরুণ}}, with \bn{ণ} at the end
\end{itemize}

\begin{tcolorbox}[breakable,colback=teal!4,colframe=teal!35!black,title={Before you move on}]
Which letter opens \textbf{\bn{অভিজ্ঞতা}}? (\textbf{\bn{অ}}.) Which vowel does it write? (\emph{ô}, the one every consonant carries.) What does \textbf{\bn{আ}} add to it? (One upright stroke.)
\end{tcolorbox}
\section[\texorpdfstring{\bn{আমার} \bn{ভালো} \bn{লাগছে} (āmār bhālo lāgchhe)}{আমার ভালো লাগছে (āmār bhālo lāgchhe)}]{\bn{আমার} \bn{ভালো} \bn{লাগছে} (āmār bhālo lāgchhe) — how I feel}

\label{lesson:BN-C272-bhalo-lagchhe}



\begin{quote}
Say \textquotedblleft{}I had a new experience in the village\textquotedblright{}. (\emph{grāme āmār ekṭā nôtun ôbhigyôtā hôlô}.) Then \textbf{\bn{ভালো} \bn{লাগা}}, \textquotedblleft{}to like\textquotedblright{}: \textquotedblleft{}good sticks to me\textquotedblright{}.
\end{quote}

\subsection*{You'll want to know: feelings now}
\begin{quote}
\textbf{\bn{আমার} \bn{এখানে} \bn{ভালো} \bn{লাগছে।}}
\end{quote}

\begin{quote}
\emph{āmār ekhāne bhālo lāgchhe.} — \textquotedblleft{}I'm enjoying it here.\textquotedblright{}
\end{quote}

Many feelings come to a person, as a thought does: the one who feels is in the \textbf{-\bn{র}} form. Going on now, they take \textbf{-\bn{ছে}}: \textbf{\bn{আমার} \bn{ভালো} \bn{লাগছে}}, \textquotedblleft{}I'm enjoying it\textquotedblright{}; \textbf{\bn{আমার} \bn{ভয়} \bn{করছে}} (\emph{āmār bhôy kôrchhe}) or \textbf{\bn{আমার} \bn{ভয়} \bn{লাগছে}} (\emph{āmār bhôy lāgchhe}), \textquotedblleft{}I'm scared\textquotedblright{}, with the same \textbf{\bn{লাগা}} as \textbf{\bn{ভালো} \bn{লাগা}}; \textbf{\bn{আমার} \bn{রাগ} \bn{হচ্ছে}} (\emph{āmār rāg hôchchhe}), \textquotedblleft{}I'm getting angry\textquotedblright{}. Words that describe a person take \textbf{\bn{আমি}}, as before: \textbf{\bn{আমি} \bn{খুশি}}, \textquotedblleft{}I'm happy\textquotedblright{}; \textbf{\bn{আমি} \bn{ক্লান্ত}}.

\subsection*{You'll want to know: feelings then}
In the past, a whole feeling takes the plain past, and a feeling that lasted takes \textbf{-\bn{ছিল}}:

\begin{quote}
\textbf{\bn{গ্রামটা} \bn{আমার} \bn{ভালো} \bn{লাগল}, \bn{কিন্তু} \bn{ঝড়ে} \bn{আমার} \bn{ভয়} \bn{করছিল।}}
\end{quote}

\begin{quote}
\emph{grāmṭā āmār bhālo lāglô, kintu jhôṛe āmār bhôy kôrchhilô.} — \textquotedblleft{}I liked the village, but I was scared in the storm.\textquotedblright{}
\end{quote}

With \textbf{\bn{হওয়া}}, a describing word becomes a change: \textbf{\bn{সে} \bn{খুশি} \bn{হল}} (\emph{she khushi hôlô}), \textquotedblleft{}she was pleased\textquotedblright{}.

\subsection*{Your turn}
Say these aloud:

\begin{itemize}
\raggedright
  \item \textquotedblleft{}I'm enjoying it here\textquotedblright{} — \emph{āmār ekhāne bhālo lāgchhe}
  \item \textquotedblleft{}I'm scared\textquotedblright{} — \emph{āmār bhôy kôrchhe}
  \item \textquotedblleft{}She was pleased\textquotedblright{} — \emph{she khushi hôlô}
\end{itemize}

\begin{tcolorbox}[breakable,colback=teal!4,colframe=teal!35!black,title={Before you move on}]
\textquotedblleft{}I'm scared\textquotedblright{}? (\textbf{\bn{আমার} \bn{ভয়} \bn{করছে}}.) \textquotedblleft{}I'm happy\textquotedblright{}? (\textbf{\bn{আমি} \bn{খুশি}}.) Which feelings take \textbf{\bn{আমার}}? (Those that come to you: \textbf{\bn{ভালো} \bn{লাগা}}, \textbf{\bn{ভয়}}, \textbf{\bn{রাগ}}.)
\end{tcolorbox}
\section[\texorpdfstring{… \bn{হতে} \bn{চাই} (… hôte chāi)}{… হতে চাই (… hôte chāi)}]{… \bn{হতে} \bn{চাই} (hôte chāi) — what I want to become}

\label{lesson:BN-C272-hote-chai}



\begin{quote}
Say \textquotedblleft{}I'm enjoying it here\textquotedblright{}. (\emph{āmār ekhāne bhālo lāgchhe}.) Then \textquotedblleft{}I have to go\textquotedblright{}. (\emph{āmāke jete hôbe}.) Its \textbf{\bn{যেতে}} is the \textbf{-\bn{তে}} form you need now.
\end{quote}

\subsection*{You'll want to know: want to}
\begin{quote}
\textbf{\bn{আমি} \bn{ডাক্তার} \bn{হতে} \bn{চাই।}}
\end{quote}

\begin{quote}
\emph{āmi ḍāktār hôte chāi.} — \textquotedblleft{}I want to be a doctor.\textquotedblright{}
\end{quote}

\textbf{\bn{চাই}} is \textbf{\bn{চাওয়া}}, \textquotedblleft{}to want\textquotedblright{}, in the present: \textquotedblleft{}I want\textquotedblright{}. With the \textbf{-\bn{তে}} form of another verb it is \textquotedblleft{}want to\textquotedblright{}: \textbf{\bn{হতে} \bn{চাই}}, \textquotedblleft{}want to be\textquotedblright{}; \textbf{\bn{দেখতে} \bn{চাই}} (\emph{dekhte chāi}), \textquotedblleft{}want to see\textquotedblright{}; \textbf{\bn{যেতে} \bn{চাই}} (\emph{jete chāi}), \textquotedblleft{}want to go\textquotedblright{}. Here the one who wants is \textbf{\bn{আমি}}, not \textbf{\bn{আমার}}: wanting is something you do.

\subsection*{You'll want to know: my wish}
For a wish, use \textbf{\bn{ইচ্ছা}}, \textquotedblleft{}a wish\textquotedblright{}, or \textbf{\bn{স্বপ্ন}}, \textquotedblleft{}a dream\textquotedblright{}, which you know, with the \textbf{-\bn{র}} form of the verb:

\begin{quote}
\textbf{\bn{আমার} \bn{পাহাড়ে} \bn{যাওয়ার} \bn{ইচ্ছা} \bn{আছে।}}
\end{quote}

\begin{quote}
\emph{āmār pāhāṛe jāoyār ichchhā āchhe.} — \textquotedblleft{}I'd like to go to the mountains.\textquotedblright{}
\end{quote}

\textbf{\bn{যাওয়ার} \bn{ইচ্ছা}}, \textquotedblleft{}a wish of going\textquotedblright{}. Said aloud, \textbf{\bn{ইচ্ছা}} is often \emph{ichchhe}, and you will see it written \textbf{\bn{ইচ্ছে}} too.

\subsection*{Your turn}
Say these aloud:

\begin{itemize}
\raggedright
  \item \textquotedblleft{}I want to see the sea\textquotedblright{} — \emph{āmi sômudrô dekhte chāi}
  \item \textquotedblleft{}She wants to be a doctor\textquotedblright{} — \emph{she ḍāktār hôte chāy}
  \item \textquotedblleft{}I'd like to go to the mountains\textquotedblright{} — \emph{āmār pāhāṛe jāoyār ichchhā āchhe}
\end{itemize}

\begin{tcolorbox}[breakable,colback=teal!4,colframe=teal!35!black,title={Before you move on}]
\textquotedblleft{}I want to be\textquotedblright{}? (\textbf{\bn{হতে} \bn{চাই}}.) \textbf{\bn{আমি}} or \textbf{\bn{আমার}} with \textbf{\bn{চাই}}? (\textbf{\bn{আমি}}.) And with \textbf{\bn{ইচ্ছা} \bn{আছে}}? (\textbf{\bn{আমার}}.)
\end{tcolorbox}
\section[(review)]{Review — feelings and hopes}

\label{lesson:BN-R272-feelings-and-hopes}



\begin{quote}
Say \textquotedblleft{}I'm scared\textquotedblright{}, then \textquotedblleft{}I want to be a doctor\textquotedblright{}. (\emph{āmār bhôy kôrchhe}; \emph{āmi ḍāktār hôte chāi}.)
\end{quote}

\subsection*{Your turn}
\begin{itemize}
\raggedright
  \item \emph{Say it:} I'm enjoying it here, because the people are kind.
  \item \emph{Say it:} Yesterday I was scared in the storm.
  \item \emph{Say it:} I had a new experience in the city.
  \item \emph{Say it:} Mira wants to be a doctor. Arun wants to see the sea.
  \item \emph{Say it:} I'd like to go to the mountains one more time.
  \item \emph{Write it:} \textbf{\bn{অভিজ্ঞতা}}, then \textbf{\bn{অরুণ}}
\end{itemize}

\begin{tcolorbox}[breakable,colback=teal!4,colframe=teal!35!black,title={Before you move on}]
\textquotedblleft{}I'm enjoying it\textquotedblright{}? (\textbf{\bn{আমার} \bn{ভালো} \bn{লাগছে}}.) \textquotedblleft{}I want to see\textquotedblright{}? (\textbf{\bn{দেখতে} \bn{চাই}}.) Which letter opens \textbf{\bn{অভিজ্ঞতা}}? (\textbf{\bn{অ}}.)
\end{tcolorbox}
